
\exo

Écrire les nombres suivants sous la forme $a^n$ ($a$ et $n$ entiers relatifs).
\[A=5^3 \times 5^8
\qpvq
B=7^{-4}\times 7^5
\qpvq
C=(-3)^{6}\times (-3)^9
\qpvq
D=\frac{4^7}{4^4}
\qpvq
E=\frac{10^{14}}{10^{20}}
\qpvq
F=\frac{9^{-5}}{9^{11}}.
\]

\exo

Parmi les nombres suivants, lesquels valent $7^5$ ?
\[\big(7^2\big)^3
\qpvq
7^2\times 7^3
\qpvq
\frac{7^2}{7^3}
\qpvq
7^8\times 7^{-3}
\qpvq
\frac{7^9}{7^4}
\qpvq
\big(7^{-1}\big)^{-5}
\qpvq
\frac{7^{12}\times 7^5}{\big(7^4\big)^3}
\qpvq
\frac{7^4}{7^{-9}\times 7^8}.
\]

\exo

\begin{enumerate}
\item Écrire les nombres suivants sous la forme d'une puissance de $3$.
\[A=27\times 9^5
\qpvq
B=\frac{81^4\times 9^{-2}}{27^{-3}}
\qpvq
C=9^{-3}\times 81^2.
\]
\item Écrire les nombres suivants sous la forme d'une puissance de $2$.
\[A=2^n\times 2^{n-1}
\qpvq
B=\frac{2^{3-n}\times 2^{n+2}}{2^{-4n}}
\qpvq
C=\big(2^{n-1}\big)^3\times 2^{-n}.
\]
\end{enumerate}

\exo

Écrire les nombres suivants sous la forme $x^n$ ($x\in\R^*$ et $n\in\Z$).
\[A=x^4\times \big(x^{-2}\big)^3\times x^5
\qpvq
B=\frac{x^7}{x^8\times x^{-5}}.\]

\exo

Écrire les nombres suivants sous la forme $a^n\times b^p$ ($a$ et $b$ dans $\R^*$, $n$ et $p$ dans $\Z$).
\[
A=\big(a^3\times b\big)^2
\qpvq
B=\left(\frac{a}{b}\right)^3\times a^{-2}
\qpvq
C=\big(a^{-5}\times b^{-2}\big)^{10}
\qpvq
D=\left(\frac{a}{b^2}\right)^{-2}\times \left(\frac{1}{a^3}\right)^{5}.
\]

\exo

\begin{enumerate}
\item Rappeler le tableau de variation de la fonction cube.
\item Sans les calculer, comparer les nombres suivants.
\vspace{-6pt}
\setlength{\columnsep}{1cm}
\setlength{\columnseprule}{0pt}
\begin{multicols}{4}
\begin{enumerate}
\item $2,5^3$\; et \; $4^3$.$\vphantom{\left(\dfrac{8}{9}\right)^3}$
\item $(-4)^3$\; et \; $(-7)^3$.$\vphantom{\left(\dfrac{8}{9}\right)^3}$
\item $3,5^3$\; et \; $27$.$\vphantom{\left(\dfrac{8}{9}\right)^3}$
\item $\left(\dfrac{5}{9}\right)^3$\; et \; $\left(\dfrac{8}{7}\right)^3$.
\end{enumerate}
\end{multicols}
\vspace{-6pt}
\end{enumerate}

\exo

Dans chacun des cas suivants, donner un encadrement de $x^3$.
\vspace{-6pt}
\setlength{\columnsep}{1cm}
\setlength{\columnseprule}{0pt}
\begin{multicols}{4}
\begin{enumerate}
\item $0\leqslant x\leqslant 2$.$\vphantom{\dfrac{3}{2}}$
\item $-1\leqslant x< 0$.$\vphantom{\dfrac{3}{2}}$
\item $-3< x< 6$.$\vphantom{\dfrac{3}{2}}$
\item $\dfrac{1}{2}\leqslant x\leqslant \dfrac{3}{2}$.
\end{enumerate}
\end{multicols}
\vspace{-6pt}

\exo

\vspace{-6pt}
\setlength{\columnsep}{1cm}
\setlength{\columnseprule}{0pt}
\begin{multicols}{2}
\raggedcolumns
On a tracé ci-contre la représentation graphique de la fonction cube.

En utilisant ce graphique, résoudre les équations et inéquations :
\vspace{-6pt}
\setlength{\columnsep}{1cm}
\setlength{\columnseprule}{0pt}
\begin{multicols}{2}
\begin{enumerate}
\item $x^3=8$.
\item $x^3>8$.
\item $x^3<8$.
\item $x^3=-5$.
\item $x^3>-5$.
\item $x^3<-5$.
\end{enumerate}
\end{multicols}
\vspace{-6pt}
\begin{center}
\def\xmin{-3}
\def\xmax{3}
\def\ymin{-30}
\def\ymax{30}
\def\xunite{1.0cm}
\def\yunite{0.07cm}
\psset{unit=\xunite, xunit=\xunite, yunit=\yunite, algebraic=true, dimen=middle, dotstyle=*, dotsize=3pt 0, linewidth=0.8pt, arrowsize=3pt 2, arrowinset=0.25, comma=true}
\begin{pspicture*}(\xmin,\ymin)(\xmax,\ymax)
%\psgrid[xunit=\xunite, yunit=\yunite, subgriddiv=0, gridlabels=0, gridcolor=lightgray](0,0)(\xmin,\ymin)(\xmax,\ymax)
\multido{\n=-3+1}{7}{\psline[linecolor=lightgray](\n,\ymin)(\n,\ymax)}
\multido{\n=-30+10}{7}{\psline[linecolor=lightgray](\xmin,\n)(\xmax,\n)}
\psaxes[labels=none,labelFontSize=\scriptstyle, xAxis=true, yAxis=true, Dx=1, Dy=10, ticksize=-2pt 0, subticks=1]{->}(0,0)(\xmin,\ymin)(\xmax,\ymax)[{\footnotesize $x$},135][{\footnotesize $y$},-45]
\psplot[linewidth=1.2pt, plotpoints=200]{\xmin}{\xmax}{x^3}
\begin{footnotesize}
\uput{3pt}[270](1,0){$1$}
\uput{3pt}[180](0,10){$10$}
\end{footnotesize}
\end{pspicture*}
\end{center}
\end{multicols}
\vspace{-6pt}

\exo

Résoudre dans $\R$ les équations et inéquations suivantes.
\vspace{-6pt}
\setlength{\columnsep}{1cm}
\setlength{\columnseprule}{0pt}
\begin{multicols}{3}
\begin{enumerate}
\item $8<x^3<27$.
\medskip
\item $x^3\geqslant 64$.
\medskip
\item $x^3>-8$.
\medskip
\item $-125<x^3<125$.

\item $4x^3>\np{4000}$.
\medskip
\item $x^3\leqslant \np{8000}$.
\medskip
\item $5x^3+2<x^3+18$.
\medskip
\item $5-2x^3\leqslant 4x^3-11$.

\item $2x^3+2<0$.
\medskip
\item $2(x-4)^3+2<0$.
\medskip
\item $2(x+5)^3-16>0$.
\medskip
\item $(x^2+1)(x-2)=-2x^2+x+3$.
\end{enumerate}
\end{multicols}
\vspace{-6pt}

\exo

Un entreprise fabrique des dés cubiques en bois, de masse volumique \SI[per-mode=symbol]{0,5}{\g\per\cm^3}. L'arête de chaque dé est comprise entre $4$ et \SI[retain-explicit-plus]{4,1}{\cm}.
\begin{enumerate}
\item Donner un encadrement du volume d'un dé.
\item En déduire un encadrement de la masse d'un dé.
\end{enumerate}

\exo

Un cône a une base circulaire de rayon $r$ (en \si{cm}), et sa hauteur $h$ est égale au triple de $r$.
\begin{enumerate}
\item Exprimer le volume de ce cône (en \si{\cm^3}) en fonction de $r$.
\item Déterminer $r$ de sorte que le volume du cône soit égal à \SI[retain-explicit-plus]{250}{\cm^3}. On donnera la valeur exacte et une valeur approchée à $0,1$ près.
\end{enumerate}

\exo

Sur la figure ci-dessous, on a tracé les courbes représentatives des fonctions $f$ et $g$ définies sur $\R$ par :
\[f(x)=-x^3+4x^2-x-7\qetq g(x)=x^2-3x-1.\]
\begin{center}
\def\xmin{-2}
\def\xmax{4}
\def\ymin{-8}
\def\ymax{10}
\def\xunite{1.5cm}
\def\yunite{0.3cm}
\psset{unit=\xunite, xunit=\xunite, yunit=\yunite, algebraic=true, dimen=middle, dotstyle=*, dotsize=3pt 0, linewidth=0.8pt, arrowsize=3pt 2, arrowinset=0.25, comma=true}
\begin{pspicture*}(\xmin,\ymin)(\xmax,\ymax)
\psgrid[xunit=\xunite, yunit=\yunite, subgriddiv=0, gridlabels=0, gridcolor=lightgray](0,0)(\xmin,\ymin)(\xmax,\ymax)
\psaxes[labelFontSize=\scriptstyle, xAxis=true, yAxis=true, Dx=1, Dy=2, ticksize=-2pt 0, subticks=1]{->}(0,0)(-1.99,-7.88)(\xmax,\ymax)[{\footnotesize $x$},135][{\footnotesize $y$},-45]
\psplot[linewidth=1.2pt, plotpoints=200]{\xmin}{\xmax}{-x^3+4*x^2-x-7}
\psplot[linewidth=1.2pt, plotpoints=200]{\xmin}{\xmax}{x^2-3*x-1}
\begin{footnotesize}
\uput{3pt}[225](3.5,-4.4){$\cC_f$}
\uput{3pt}[135](3.5,0.75){$\cC_g$}
\end{footnotesize}
\end{pspicture*}
\end{center}
L'objectif de cet exercice est de résoudre l'inéquation $f(x)<g(x)$.
\begin{enumerate}
\item Résoudre graphiquement l'inéquation $f(x)<g(x)$ (on utilisera des valeurs approchées à $0,1$ près).
\item On note $h$ la fonction définie sur $\R$ par :
\[h(x)=f(x)-g(x).\]
\begin{enumerate}
\item Vérifier que, pour tout réel $x$ : \quad $h(x)=-x^3+3x^2+2x-6$.
\item Montrer que, pour tout réel $x$ : \quad $h(x)=(x^2-2)(3-x)$.
\item Résoudre l'inéquation $h(x)<0$, et répondre au problème posé en introduction.
\end{enumerate}
\end{enumerate}

